\bookchapter{Ordinal Selection and the Delegated Exact
Engine}\label{viii-ch07-ordinal-selection-and-the-delegated-exact-engine}

The latest global selector described by the retained SAMA material is
Q2. Its job is to order an admissible catalog using source-visible
features, while exact arithmetic remains delegated. This chapter records
its frozen comparisons without turning a learned ordering into a
physical law or an unrestricted complexity claim.

The account below follows the chapter source as retained at the fixed 8
September 2026 SAMA revision. Engine names, installed interfaces and
campaign status are historical source statements. They do not describe a
freshly inspected October deployment. Later source amendments are
distinguished where available; no later engine or campaign is inferred.

The September SAMA selector chapter cites H000531, H000551 and H000552
for Q2. The accessible August public SAM Research Project live page
instead records V21R, and its named Q2 candidate and promotion artifact
paths were unavailable during preparation. This chapter therefore
preserves the later manuscript account with that retrieval limitation;
it does not independently certify a live pointer transition.

\textbf{Source-status qualification.} The operational and revision
statements in this chapter describe the retained source documents. They
are not an instruction to execute, deploy, resume, train, promote or
change a service, and they have not been independently verified as
present live state for this edition.

\subsection{Overview}\label{viii-ch07-overview}

Exact computation can return several admissible receipts for one query.
The arithmetic substrate certifies what every receipt is, while the
selection plane asks a different question: which admissible receipt has
the best exact return quality when the selector is permitted to see only
source geometry? \texttt{SLCQ2-RZ} is the frozen current answer to that
ordinal question.

The target quality is not a learned label invented after execution. For
each candidate receipt it is the frozen reciprocal-quotiented
lexicographic vector

\[
Q=(d_{direction},r_{inverse},r_{W8},r_{D81}),
\]

where lower is better and equal vectors remain a complete target tie. Q2
does not read any component used to construct that target. It exposes 70
exact integer features computed from axis, modulus, shell, direction and
two typed domain labels. It excludes inverse inputs, inverse directions,
returned endpoints, D81 checkpoints, target tiers, custody identities,
hashes, execution order and timing.

The data surface contains 5,120 inherited exact receipts, 1,024
canonical queries and 512 reciprocal/replicate-safe groups. Each axis is
frozen into 80 TRAIN, 24 HOLDOUT and 24 CONTROL groups. TRAIN supplies
640 queries; HOLDOUT and CONTROL each expose a complete 192-query
catalog. Three exact-rational architectures at three predeclared ridge
values form nine fits. HOLDOUT fixes the winner before CONTROL is
opened:

\nolinkurl{Q2-HYBRID_LISTWISE_ADJACENT_ORDINAL_RIDGE-RIDGE-1D1000}.

Against frozen Q1 on the same fresh surfaces, Q2 raises weighted
strict-pair accuracy by exactly \texttt{63200/594823} on HOLDOUT and
\texttt{83595/773858} on CONTROL, while recovering every exact best set,
creating no false singleton, making no abstention and producing no
strict score tie. The installed classification is exactly: \textbf{The
test result suggests strong contact with the concept.}

The architecture is deliberately layered. Q2 owns the current ordinal
selection plane. Frozen \texttt{SLCV21R} remains noncurrent as a
revision but remains the byte-identical delegated arithmetic and
application base. Q1 is an immutable fallback. Neither selection nor
promotion assigns physical scale, absorbs ATOM3D, changes Exact Write
semantics or authorizes application work.

\subsection{Opening question and conceptual
picture}\label{viii-ch07-opening-question-and-conceptual-picture}

The governing question is:

\begin{quote}
Can exact source-visible event geometry order the complete admissible
receipt catalog more faithfully than Q1 without reading any executed
value that defines the answer?
\end{quote}

The typed route is

\begin{itemize}
\tightlist
\item
  frozen SLCV21R exact query and candidate receipts
\item
  reciprocal/replicate-safe grouping
\item
  target-only exact return-quality vector
\item
  70-field source-visible feature map
\item
  exact listwise and adjacent-ordinal fitting on TRAIN
\item
  winner lock on complete HOLDOUT
\item
  one comparison on sealed complete CONTROL
\item
  frozen Q2 selection service
\item
  SLCQ2-RZ current pointer
\end{itemize}

Three separations carry the meaning of the result:

\begin{enumerate}
\def\labelenumi{\arabic{enumi}.}
\tightlist
\item
  the exact substrate determines admissibility and computes receipts;
\item
  the target oracle evaluates return quality but is hidden from
  selection;
\item
  the selector orders candidates from permitted source geometry alone.
\end{enumerate}

Collapsing these surfaces would make the question circular. Keeping them
separate lets the result say something precise: the visible structural
fields contain enough ordinal information for the selected exact model
to recover the complete best set and nearly all strict pair orderings on
two held-out complete catalogs.

\subsection{Recorded authority and architectural
planes}\label{viii-ch07-recorded-authority-and-architectural-planes}

\subsubsection{Sole source-recorded default
revision}\label{viii-ch07-sole-source-recorded-default-revision}

\texttt{SLCQ2-RZ} has status
\nolinkurl{FROZEN_COMPLETE_STOPPED_PROMOTED_CURRENT} and is the sole
source-recorded default SLC pointer. Older v1.x, V20R and
\texttt{SLCV21R} revisions retain their artifacts and scoped results;
they are not alternative current pointers. The older Volume III
technical spine therefore supplies lineage, not present-tense pointer
authority.

Q1 remains frozen and immutable as the fallback selector. A fallback is
not a second default: it is a preserved predecessor available under the
declared service boundary.

\subsubsection{Selection is not
arithmetic}\label{viii-ch07-selection-is-not-arithmetic}

For a query \(q\), let \(\mathcal C(q)\) be the finite set of receipts
certified admissible by the exact base. Q2 supplies a score

\[
s_q(c)=w^T\phi(q,c),\qquad c\in\mathcal C(q),
\]

with exact rational coefficient vector \(w\) and exact integer feature
vector \(\phi(q,c)\in\mathbb Z^{70}\). It returns all candidates
attaining the maximum permitted score, subject to the ambiguity policy
described below.

The delegated base still owns DOS construction, exact write, reciprocal
return, checkpoints and application execution. Installing a new \(w\)
changes the ordering of already-admissible candidates; it does not
change what the candidate means or how its exact receipt is computed.

\subsubsection{Four distinct planes}\label{viii-ch07-four-distinct-planes}

\begin{itemize}
\tightlist
\item
  \textbf{exact substrate.} Owner in the installed architecture: frozen
  \texttt{SLCV21R} services; Output: admissible exact receipts and
  arithmetic.
\item
  \textbf{quality target.} Owner in the installed architecture: frozen
  reciprocal-quotiented oracle; Output: exact lexicographic target
  vectors.
\item
  \textbf{ordinal selector.} Owner in the installed architecture:
  current \texttt{SLCQ2-RZ} Q2 model; Output: score order and complete
  top-score set.
\item
  \textbf{application.} Owner in the installed architecture: separately
  declared adapter/model; Output: application-specific interpretation or
  execution.
\end{itemize}

An ATOM3D model, RH adapter or Mersenne scheduler may call the service
only through its own typed contract. None becomes part of Q2 merely by
using its ordering.

\subsection{Target construction and blindness
contract}\label{viii-ch07-target-construction-and-blindness-contract}

\subsubsection{Exact target order}\label{viii-ch07-exact-target-order}

The target for candidate \(c\) is

\[
Q_q(c)=
\bigl(
d_{direction}(q,c),
r_{inverse}(q,c),
r_{W8}(q,c),
r_{D81}(q,c)
\bigr).
\]

Vectors compare lexicographically, with smaller values preferred. The
first coordinate at which two candidates differ decides their strict
order. If no coordinate differs, both belong to the same target tier. In
particular, a custody identifier never converts an exact target tie into
a singleton.

The complete best set is

\[
B_q=\{c\in\mathcal C(q):Q_q(c)=\min_{c'\in\mathcal C(q)}Q_q(c')\}.
\]

Recovery is exact only when the returned set equals \(B_q\), not when it
merely contains one member of \(B_q\).

\subsubsection{Permitted source fields}\label{viii-ch07-permitted-source-fields}

Q2 may read only:

\begin{itemize}
\tightlist
\item
  axis
\item
  modulus
\item
  left\_shell, right\_shell
\item
  left\_direction, right\_direction
\item
  route\_domain
\item
  physical\_relation\_domain
\end{itemize}

From these it computes the fixed 70-field map. The two domain labels are
constant on this source; they serve as a typed zero-variation
interaction check, not as an undeclared ranking signal.

\subsubsection{Forbidden target and custody
fields}\label{viii-ch07-forbidden-target-and-custody-fields}

The selector may not read:

\begin{itemize}
\tightlist
\item
  inverse input lanes or inverse directions
\item
  returned endpoint lanes
\item
  D81 checkpoint lanes
\item
  quality vector or quality tier
\item
  receipt index or pair identity
\item
  group hash or source hash
\item
  execution sequence
\item
  timing
\end{itemize}

This prohibition is load-bearing. Inverse, endpoint and checkpoint
values construct the oracle and would leak the answer. Custody fields
could create memorized ordering without structural meaning. Timings
could import hardware and execution-order artifacts into a mathematical
selector.

\subsection{Fresh grouped custody}\label{viii-ch07-fresh-grouped-custody}

\subsubsection{Why the unit is a physical
group}\label{viii-ch07-why-the-unit-is-a-physical-group}

Individual receipt splitting is a wrong control because reciprocal
directions and replicated views of one physical relation can then appear
on opposite sides of a split. A model could learn one view and be scored
on its mate. Q2 hashes a fresh namespace together with axis, modulus,
unordered shell endpoints and the two source domain labels. All
reciprocal directions and replicates sharing that identity remain in one
group.

The frozen census is

\begin{itemize}
\tightlist
\item
  \textbf{inherited exact receipts.} Exact count: 5,120.
\item
  \textbf{canonical queries.} Exact count: 1,024.
\item
  \textbf{reciprocal/replicate-safe groups.} Exact count: 512.
\item
  \textbf{reciprocal canonical queries per group.} Exact count: 2.
\item
  \textbf{observed group multiplicities.} Exact count: 4, 10, 16.
\end{itemize}

The nonuniform multiplicities matter. A group contributes according to
its actual candidate roster; assuming one uniform multiplicity changes
the catalog and its weights.

\subsubsection{Split derivation}\label{viii-ch07-split-derivation}

Within each axis, group hashes are sorted and frozen as

\[
80\ \text{TRAIN}+24\ \text{HOLDOUT}+24\ \text{CONTROL}=128
\]

groups. Across the four axes this gives the 512-group total. The query
census is

\[
4\cdot80\cdot2=640\quad\text{TRAIN queries},
\]

\[
4\cdot24\cdot2=192
\]

queries on each of HOLDOUT and CONTROL.

TRAIN uses the frozen harder tier rule and contains 18,654 group
references and 182,688 receipt exposures. HOLDOUT and CONTROL retain
every split-local candidate group, yielding complete 192-query catalogs.
Selection therefore cannot improve a metric by silently dropping
difficult candidates.

\subsection{Deriving the 70-feature
representation}\label{viii-ch07-deriving-the-70-feature-representation}

\subsubsection{Direct and reciprocal
alignment}\label{viii-ch07-direct-and-reciprocal-alignment}

Let the centered source shell endpoints of a query and candidate be
exact integers. Q2 forms both alignments:

\begin{itemize}
\tightlist
\item
  direct: query left -\textgreater{} candidate left, query right
  -\textgreater{} candidate right
\item
  swapped: query left -\textgreater{} candidate right, query right
  -\textgreater{} candidate left
\end{itemize}

For each alignment it retains signed residuals before taking absolute or
squared summaries. Schematically, with query endpoints ((a,b)) and
candidate endpoints ((c,d)),

\[
\delta_{direct}=(c-a,d-b),
\qquad
\delta_{swap}=(d-a,c-b).
\]

The map includes the signed components, absolute components, squares,
alignment sums and best/other residual summaries. Because both
constructions are ordered by a source-visible canonical key before
exposure, a whole-lane swap exchanges internal work without changing the
final visible vector.

\subsubsection{Midpoint, span, direction and
parity}\label{viii-ch07-midpoint-span-direction-and-parity}

For endpoints ((a,b)), define signed midpoint numerator and span

\[
m=a+b,
\qquad
\ell=b-a.
\]

Using the numerator avoids a floating-point half. Candidate-minus-query
relations for \(m\) and \(\ell\) retain orientation and separation that
coarse shell bins erased. Individual direction matches and endpoint
parity relations are then appended rather than compressed to a single
flag.

\subsubsection{Two-adic profiles and typed
interactions}\label{viii-ch07-two-adic-profiles-and-typed-interactions}

For a nonzero integer \(n\), \(v_2(n)\) is the largest \(e\ge0\) for
which \(2^e\mid n\). Q2 includes four direct/swapped alignment
valuations and two span-relation valuations, with the frozen convention
for zero. These features expose dyadic shell structure without reading
any returned value.

Route-domain and physical-relation-domain equality fields and their
declared interactions preserve the type boundary. Fourteen bounded
quadratic cross terms add fixed pair interactions; each is clipped at
exact magnitude 16,777,215. The clip is part of the pre-execution
feature contract and not a post-result adjustment.

\subsubsection{Why Q2 succeeds where the 20-field map
collided}\label{viii-ch07-why-q2-succeeds-where-the-20-field-map-collided}

Q1's coarser 20-feature map has strict-pair feature-collision rates
2.3255657565\% on HOLDOUT and 2.1898074324\% on CONTROL. When two
candidates with different target order have the same visible vector, any
deterministic linear score must tie them. Q2's signed, unbinned and
reciprocal-invariant map has zero strict feature collisions on both
surfaces. This removes that representational obstruction before model
fitting; it does not guarantee a correct order by itself.

\subsection{Exact ordinal fitting}\label{viii-ch07-exact-ordinal-fitting}

\subsubsection{Three predeclared
architectures}\label{viii-ch07-three-predeclared-architectures}

The bounded roster contains:

\begin{enumerate}
\def\labelenumi{\arabic{enumi}.}
\tightlist
\item
  \nolinkurl{EXPANDED_QUERY_CENTERED_LISTWISE_RIDGE}, which fits
  query-centered tier grades over the harder TRAIN catalog;
\item
  \nolinkurl{HYBRID_LISTWISE_ADJACENT_ORDINAL_RIDGE}, which appends
  exact differences for adjacent strict tiers with frozen weight four;
  and
\item
  \nolinkurl{LEXICOGRAPHIC_CONDITIONAL_MULTI_HEAD_RIDGE}, which learns
  the first target coordinate globally and later coordinates only inside
  equality classes of the preceding coordinates.
\end{enumerate}

Each architecture is fitted at exact ridge values

\[
\lambda\in\left\{\frac1{1000},\frac1{100},\frac1{10}\right\},
\]

giving \(3\times3=9\) fits. There is no floating-point training and no
full-scale hyperparameter sweep.

\subsubsection{Listwise and adjacent
observations}\label{viii-ch07-listwise-and-adjacent-observations}

For each training query, a target tier grade is centered within its
complete candidate list. The listwise design therefore learns relative
order without letting different query offsets dominate.

For adjacent strict tiers, choose \(c_+\) from the better tier and
\(c_-\) from the next worse tier. The hybrid supplies the exact
difference row

\[
\Delta\phi=\phi(q,c_+)-\phi(q,c_-),
\]

with a positive ordinal target and frozen weight four. A fitted vector
should therefore satisfy

\[
w^T\Delta\phi>0
\]

for that adjacent relation. Chaining adjacent relations transports the
order through the tier sequence while the listwise rows retain the
complete query geometry.

With exact design matrix \(X\), target \(y\), diagonal frozen
observation weights \(W\) and ridge \(\lambda\), the rational normal
system is

\[
(X^TWX+\lambda I)w=X^TWy.
\]

All inputs and outputs remain \texttt{Fraction} values. Exact fitting
makes score equality a mathematical event rather than a floating
tolerance choice.

\subsubsection{Winner and tie policy}\label{viii-ch07-winner-and-tie-policy}

The three hybrid fits have identical complete-catalog ranking metrics.
Only after mathematical metrics and parameter count tie does the frozen
key order select the lower regularization identifier. The resulting
winner is

\nolinkurl{Q2-HYBRID_LISTWISE_ADJACENT_ORDINAL_RIDGE-RIDGE-1D1000}.

At service time, a score tie spanning distinct target vectors triggers
abstention. A score tie wholly inside one exact target vector returns
the complete tied best set. The promoted wrapper names a top-score
multiplicity explicitly as \texttt{AMBIGUOUS\_TOP\_SCORE}; it does not
choose by receipt ID.

\subsection{HOLDOUT selection and sealed
CONTROL}\label{viii-ch07-holdout-selection-and-sealed-control}

\subsubsection{Predeclared material
comparator}\label{viii-ch07-predeclared-material-comparator}

Before execution, success required all three conditions:

\begin{itemize}
\tightlist
\item
  at least one percentage point of weighted strict-pair gain over frozen
  Q1 on HOLDOUT;
\item
  positive gain on CONTROL; and
\item
  perfect exact-best-set recovery on both complete surfaces.
\end{itemize}

HOLDOUT selected and locked the model. CONTROL was opened once afterward
for the comparison to the frozen Q1 baseline.

\subsubsection{Exact results}\label{viii-ch07-exact-results}

\begin{itemize}
\tightlist
\item
  \textbf{HOLDOUT.} Frozen Q1: \texttt{530826/594823}; Q2 winner:
  \texttt{594026/594823}; Exact delta: \texttt{63200/594823}; Decimal
  point gain: +10.625009.
\item
  \textbf{CONTROL.} Frozen Q1: \texttt{689509/773858}; Q2 winner:
  \texttt{386552/386929}; Exact delta: \texttt{83595/773858}; Decimal
  point gain: +10.802369.
\end{itemize}

The CONTROL Q2 fraction has denominator half the Q1 display denominator,
so the like-denominator comparison is

\[
\frac{386552}{386929}
=\frac{773104}{773858},
\]

and

\[
\frac{773104-689509}{773858}
=\frac{83595}{773858}.
\]

Q2 retains 6,376 wrong weighted strict pairs on HOLDOUT and 6,032 on
CONTROL. Those residual errors remain in the record. At the service
boundary, however, best-tier accuracy and exact-best-set recovery are
both \(1/1\) on both splits, with zero false singletons, zero
abstentions and zero strict score ties.

The authorized classification is:

\textbf{The test result suggests strong contact with the concept.}

\subsection{Deviation, correction and retest
chain}\label{viii-ch07-deviation-correction-and-retest-chain}

\subsubsection{Candidate-construction
failures}\label{viii-ch07-candidate-construction-failures}

Five preserved attempts materially shaped the frozen candidate:

\begin{itemize}
\tightlist
\item
  \textbf{001.} Failure: a cross-axis unit fixture mixed an
  axis-specific assumption into the test; Correction: construct the
  fixture inside its declared axis/type; Retest consequence: boundary
  tests reached the intended feature question.
\item
  \textbf{002.} Failure: code assumed uniform group multiplicity;
  Correction: retain the frozen observed roster \texttt{\{4,10,16\}};
  Retest consequence: complete catalogs and exposure weights
  reconstructed exactly.
\item
  \textbf{003.} Failure: reciprocal orientation changed the visible
  feature vector; Correction: canonicalize complete direct/swapped
  blocks by source-visible keys; Retest consequence: whole-lane swap
  became invariant.
\item
  \textbf{004.} Failure: a serializer could not carry the exact fraction
  surface; Correction: encode exact numerators/denominators through the
  declared receipt form; Retest consequence: checkpoint and replay
  identity remained exact.
\item
  \textbf{005.} Failure: the Q1 collision diagnostic read the wrong
  comparison source; Correction: recompute that diagnostic from the
  frozen Q1 comparator on Q2 surfaces; Retest consequence: the corrected
  Q1 collision rates changed; Q2 model selection and CONTROL result did
  not.
\end{itemize}

The first four attempts stopped before model selection or CONTROL
access. The fifth corrected a predecessor diagnostic only. Keeping these
attempts makes clear which faults affected construction and which did
not affect the frozen winner.

\subsubsection{Promotion-wrapper
corrections}\label{viii-ch07-promotion-wrapper-corrections}

Promotion also preserved two distinct corrections:

\begin{itemize}
\tightlist
\item
  Attempt 001 expected an \texttt{assertions\_passed} field, while the
  independent receipt records \texttt{assertion\_count=165}. The wrapper
  was corrected to read the actual schema before installation continued.
\item
  Attempt 002 installed an executable resolution correctly, but the
  pointer inherited explanatory usage language and stale promotion
  fields. Pointer metadata was corrected, producing current pointer
  SHA-256
  \nolinkurl{751f6e6a6dbe01cfdc335cf303824887148c9eb657ded691264e37c6f61fa318},
  and the promotion surface was revalidated.
\end{itemize}

Neither correction retrained Q2, reopened Q1 or mined CONTROL.

\subsubsection{Wrong controls retained as
lessons}\label{viii-ch07-wrong-controls-retained-as-lessons}

The following alternatives would answer a different question and are not
the installed controls:

\begin{itemize}
\tightlist
\item
  receipt-level random splitting, because it leaks reciprocal or
  replicated physical relations;
\item
  incomplete HOLDOUT or CONTROL candidate rosters, because they can hide
  misranking and false singleton behavior;
\item
  target-visible fields, because they make ordinal prediction circular;
\item
  custody-ID tie breaking, because it converts exact ambiguity into an
  arbitrary order;
\item
  floating score tolerances, because exact equality would depend on
  runtime;
\item
  opening CONTROL during model choice, because CONTROL would cease to be
  the postselection surface.
\end{itemize}

\subsection{Freeze and promotion
process}\label{viii-ch07-freeze-and-promotion-process}

The candidate freeze passed 8/8 release checks, 14/14 reconstruction
checks, 165/165 independent assertions and, after the diagnostic
correction, 8/8 release checks again. The promotion wrapper then passed:

\begin{itemize}
\tightlist
\item
  \textbf{pre-install.} Result: 22/22.
\item
  \textbf{post-install.} Result: 12/12.
\item
  \textbf{independent promotion validation.} Result: 20/20.
\item
  \textbf{final release.} Result: 6/6.
\end{itemize}

Promotion changed the live selection pointer through H000551; H000552
records the pointer-metadata correction. It preserved the candidate, Q1
and delegated base artifacts rather than rewriting them.

The promotion performed no retraining, CONTROL mining, physical
calibration, ATOM3D specialization, remote transfer, prime work,
publication or deployment. These are not missing steps in Q2. They are
separate campaigns requiring their own contracts and authority.

\subsection{Established result and
boundary}\label{viii-ch07-established-result-and-boundary}

This document establishes the following typed statement:

\begin{quote}
On the frozen Q2 grouped catalogs, the 70-feature exact hybrid
listwise-plus-adjacent ordinal selector materially improves strict-pair
ordering over frozen Q1 and recovers every complete exact best set. Its
promotion makes \texttt{SLCQ2-RZ} the sole source-recorded default SLC
selection revision, while exact arithmetic remains delegated to frozen
\texttt{SLCV21R} and Q1 remains the immutable fallback.
\end{quote}

It does not establish a physical energy scale, a new DOS, a change to
Exact Write, a global optimizer for every application, native GPU
credit, ATOM3D specialization, an RH conclusion or a Mersenne status.
The next chapter uses Q2 only as a typed application dependency and
keeps its Mersenne scheduler and certificates separate.

\subsection{Sources and evidential
scope}\label{viii-ch07-sources-and-evidential-scope}

This chapter is assembled from the following fixed source manuscripts.
Source-reported campaign outcomes retain their original scope; no new
campaign was executed for this edition.

\begin{itemize}
\tightlist
\item
  \href{https://github.com/iwtbotiwtwot/SAMA/blob/4c2b5dc1649da5900c90b90d0ce56c223dbb7717/vol_iii/02-slc/slcq2-rz-ordinal-selection-engine.md}{Slcq2
  rz ordinal selection engine}.
\end{itemize}
